\documentclass[12pt]{article}
\usepackage{amsmath,amssymb,amsfonts,amsthm}
\usepackage[margin=1in]{geometry}

\begin{document}

\begin{center}
{\Large\bfseries Chapter 4, Problem 24}\\[6pt]
Byron \& Fuller, \textit{Mathematics of Classical and Quantum Physics}
\end{center}

\bigskip

\textbf{Problem.} Define an operator $P_i$ by
\[
P_i = \sum_j |x_j^{(i)}\rangle\langle x_j^{(i)}|,
\]
where $i = 1, \ldots, p$, $\mu_i$ is the degree of degeneracy of the $i$th eigenvalue of some operator $A_0$, and $x_j^{(i)}$ are the corresponding eigenvectors. Let $\epsilon A_1$ be the perturbation. Show that if
\[
\epsilon A_1 = \sum_{i,k} P_i\, \epsilon A_1\, P_k, \quad A_1 = \sum_i P_i A_1 P_i,
\]
then it is possible to solve for the eigenfunctions and eigenvalues of $A_0$ in terms of the eigenfunctions and eigenvalues of $A_0$. If the degeneracy of $A_0$ is completely lifted in first order by $A_1$, then show that we can do ordinary nondegenerate perturbation theory on $A_1'$. Show that the first-order corrections to the eigenvectors of $A_0$ are all zero. Make any comments that seem relevant about the relative sizes of terms in the first few orders of perturbation theory.

\bigskip
\textbf{Solution.}

The projection operators $P_i$ project onto the degenerate subspace belonging to eigenvalue $\lambda_i^{(0)}$ of $A_0$. They satisfy:
\[
P_i P_j = \delta_{ij} P_i, \quad \sum_i P_i = I, \quad A_0 P_i = \lambda_i^{(0)} P_i.
\]

\textbf{Condition $A_1 = \sum_i P_i A_1 P_i$:}

This means $A_1$ has no matrix elements between different degenerate subspaces:
\[
P_i A_1 P_k = 0 \quad \text{for } i \neq k.
\]

In other words, $A_1$ is block-diagonal with respect to the decomposition of the Hilbert space into eigenspaces of $A_0$.

\textbf{Solving within each subspace:}

Within the $i$th degenerate subspace, we diagonalize the $\mu_i \times \mu_i$ matrix:
\[
\langle x_j^{(i)}|A_1|x_k^{(i)}\rangle.
\]
The eigenvalues $\lambda_{j}^{(1,i)}$ give the first-order corrections, and the eigenvectors give the correct zeroth-order states.

\textbf{If degeneracy is completely lifted:}

If all the eigenvalues $\lambda_n^{(0)} + \epsilon\lambda_n^{(1)}$ are distinct after first-order correction, then from second order onward we can apply nondegenerate perturbation theory. The ``effective'' unperturbed operator is $A_0' = A_0 + \epsilon P_i A_1 P_i$ (within each subspace), and the remaining perturbation (off-diagonal terms connecting different subspaces, which are of order $\epsilon$ but connect states separated by zeroth-order energy gaps) can be treated perturbatively.

\textbf{First-order corrections to eigenvectors are zero:}

The first-order eigenvector correction in nondegenerate perturbation theory is:
\[
|x_n^{(1)}\rangle = \sum_{m \neq n} \frac{\langle x_m^{(0)}|A_1|x_n^{(0)}\rangle}{\lambda_n^{(0)} - \lambda_m^{(0)}} |x_m^{(0)}\rangle.
\]

Since the condition $A_1 = \sum_i P_i A_1 P_i$ means $\langle x_m^{(0)}|A_1|x_n^{(0)}\rangle = 0$ whenever $x_m^{(0)}$ and $x_n^{(0)}$ belong to different eigenspaces of $A_0$ (i.e., $\lambda_m^{(0)} \neq \lambda_n^{(0)}$), all terms in the sum vanish. Therefore:
\[
|x_n^{(1)}\rangle = 0.
\]

\textbf{Comments on relative sizes:}

Since the first-order eigenvector correction vanishes, the leading correction to the eigenvectors is of order $\epsilon^2$ rather than $\epsilon$. The perturbation series for eigenvalues has the structure $\lambda = \lambda^{(0)} + \epsilon\lambda^{(1)} + \epsilon^3\lambda^{(3)} + \cdots$ (the second-order eigenvalue correction also vanishes because it depends on the first-order eigenvector correction). The effective perturbation expansion converges faster than the standard case. \qed

\end{document}
